\section{Conclusion}
\label{sec:conclusion}

Bornholt et al.'s $\Utype\langle T \rangle$ improved how programs handle
noisy measurements by lifting random variables into the type system and
discharging conditionals through sequential testing. The abstraction
assumes the measurement exists. Many real workloads do not satisfy that
assumption: clinical reports go missing, sensors drop readings, GPS fixes
vanish in tunnels. Treating absent data as a wide distribution leaks
fabricated numbers; wrapping it as $\texttt{Option<Uncertain<T>>}$
collapses a probability into a Boolean and reintroduces the uncertainty
bugs the original type was designed to eliminate.

$\Mtype\langle T \rangle$ closes this gap with a small, principled
extension. The type is a product of two independent
$\Utype$-valued random variables, one for presence and one for value;
arithmetic propagates both axes, with presence combining under AND; a
single new operator, $\textsc{lift}$, runs SPRT on the presence axis and
returns a plain $\Utype\langle T \rangle$ only when the evidence
supports the presence hypothesis at the caller's chosen confidence. The
implementation reuses the existing $\Utype$ machinery and adds roughly
two hundred lines of Rust across two files. Per-sample overhead is a
constant factor between $1.2\times$ and $2.2\times$; SPRT-gated lifting
costs a few hundred nanoseconds in decisive regimes and a few
microseconds when the operand's presence probability sits inside the
indifference region.

The construction does not claim mathematical novelty. The factoring of
presence from value appears in probabilistic databases, in Kalman
filtering with intermittent observations, and in probabilistic
programming idioms. The contribution is making the factoring available
as a value-level type in a general-purpose programming language, with
operator-level support and a statistically-gated bridge back to the
unfactored type. The principal direction for future work is dropping the
independence assumption to support missing-not-at-random workloads
through a joint presence-value distribution; supporting more numeric
specialisations beyond $f64$ and $\Bool$; and lifting compound Boolean
expressions on $\Mtype\langle \Bool \rangle$ to first-class operators.
A complementary direction is formal verification of programs that use
$\Mtype\langle T \rangle$ using one of the recent relational
probabilistic program logics.

\newpage